\section{Reproducibility specification}
\label{app:repro}

This appendix documents the hook paths, consumed-token mask, projection and probe, control constructions, cohort roles, and protocol record. Every statement below comes from the released adapters (\texttt{src/cftransfer/adapters/}), the registered protocol file, and the per-block artifacts.

\subsection{Consumed loci per family}
\label{app:repro-loci}

Table~\ref{tab:cf-loci} gives the primary and connector locus module paths exactly as the adapters name them, together with the rule for identifying the consumed final block. Each adapter resolves the block index at load time from the loaded checkpoint's own configuration. Each block's preflight file records the resolved path, hidden dimension, consumer, and any branch that bypasses the locus. Table~\ref{tab:cf-models} gives the consumed-token count per image at the primary locus for every checkpoint.

\begin{table}[h]
\centering
\scriptsize
\setlength{\tabcolsep}{4pt}
\caption{\textbf{Hook paths per checkpoint family.} Module paths as the adapters name them under the loaded model. $n$ is the length of the block list of that checkpoint's tower; the adapter resolves it at load time and the block's preflight file records the resolved path.}
\label{tab:cf-loci}
\begin{tabular}{>{\raggedright\arraybackslash}p{0.15\textwidth}>{\raggedright\arraybackslash}p{0.25\textwidth}>{\raggedright\arraybackslash}p{0.20\textwidth}>{\raggedright\arraybackslash}p{0.28\textwidth}}
\toprule
Family (adapter) & Primary locus & Connector locus & Consumed final block identified by \\
\midrule
Qwen2.5-VL, Qwen3-VL, Lingshu (\texttt{qwen}) & \texttt{model.visual.blocks.}$(n{-}1)$ & \texttt{model.visual.merger} & The merger reads the last entry of \texttt{model.visual.blocks}: $n=32$ for Qwen2.5-VL and Lingshu, $24$ for Qwen3-VL 4B, $27$ for Qwen3-VL 8B and 32B. The Qwen3-VL DeepStack mergers read earlier blocks and are not hooked. \\
\addlinespace[2pt]
Gemma~3, MedGemma (\texttt{gemma}) & \texttt{model.vision\_tower.}\allowbreak\texttt{encoder.layers.26} & \texttt{model.}\allowbreak\texttt{multi\_modal\_projector} & The projector reads the last SigLIP encoder layer, 26 of 27. \texttt{post\_layernorm} sits between that layer and the projector. \\
\addlinespace[2pt]
InternVL3.5 (\texttt{internvl}) & \texttt{model.vision\_tower.}\allowbreak\texttt{encoder.layer.}$(n{-}1)$ & \texttt{model.}\allowbreak\texttt{multi\_modal\_projector} & The projector reads the last InternViT encoder layer, $n=24$ at 8B and 14B and $n=45$ at 38B, after \texttt{vision\_tower.layernorm}, which is an identity under the shipped \texttt{use\_mean\_pooling} configuration. \\
\addlinespace[2pt]
LLaVA-1.5, LLaVA-Med (\texttt{llava}) & \texttt{model.vision\_tower.}\allowbreak\texttt{encoder.layers.22} & \texttt{model.}\allowbreak\texttt{multi\_modal\_projector} & The checkpoint's \texttt{vision\_feature\_layer} is $-2$, so the projector reads CLIP encoder layer 22 of 24 and the final layer is computed and discarded. \\
\addlinespace[2pt]
Llama~3.2 Vision (\texttt{mllama}) & \texttt{model.vision\_model.}\allowbreak\texttt{global\_transformer.layers.7} & \texttt{model.}\allowbreak\texttt{multi\_modal\_projector} & The projector input concatenates the last global-transformer layer, 7 of 8, with the outputs of local layers 3, 7, 15, 23 and 30; those five bypass the primary locus. \\
\bottomrule
\end{tabular}
\end{table}
\FloatBarrier

\subsection{The consumed-token mask}
\label{app:repro-mask}

The write of Equation~\ref{eq:intervention} changes only the entries of the hooked tensor corresponding to token positions the consumer reads at that locus. The adapter derives these positions for each batch from the processor's own outputs. Two tensor shapes occur. The Qwen family passes the merger a flat $(N,D)$ tensor in which every image occupies one contiguous slice. The adapter computes the slice bounds from \texttt{image\_grid\_thw} (the connector slice is a quarter of the length, one merged token per $2\times2$ patch group). Every other family passes the consumer a batched $(B,T,D)$ tensor. The adapter supplies one boolean mask per batch element:

\begin{itemize}[leftmargin=*,itemsep=1pt,topsep=2pt]
\item Gemma~3 and MedGemma: all 4{,}096 patch rows; the tower produces no CLS row. The mask length is checked against the count of \texttt{<image\_soft\_token>} placeholders in the language-model input.
\item InternVL3.5: index 0 is the CLS row, which the consumer drops, so it is excluded; the remaining rows of the single $448\times448$ tile are written. The mask length is checked against the count of \texttt{<IMG\_CONTEXT>} placeholders.
\item LLaVA-1.5 and LLaVA-Med: index 0 is the CLS row, which the checkpoint's \texttt{vision\_}\allowbreak\texttt{feature\_}\allowbreak\texttt{select\_}\allowbreak\texttt{strategy} drops under its \texttt{default} setting, so it is excluded; rows 1 to 576 are written.
\item Llama~3.2 Vision: the 1{,}601 rows (CLS and 1{,}600 patches) of the one real $560\times560$ tile. The seven zero-padding rows of each tile and the three padding tiles are excluded, and the mask is checked against \texttt{aspect\_ratio\_mask} and the last position of \texttt{cross\_attention\_mask}.
\end{itemize}

The hook raises an error rather than falling back when a flat layout does not exactly cover the hooked tensor, a mask length differs from the token-axis length, a batch element has no consumed token, or the hooked module fires more than once in one forward pass.

Preflight check (C) verifies the mask. It captures the locus tensor on a clean pass and a pass steered at $\alpha=+0.25$ along the first concept direction. It reports the largest absolute difference outside the mask and the largest absolute changes in the consumer's output and answer logits. Across all \cfBlocks{} blocks, the change outside the consumed tokens is $0$ at both loci, while both the consumer's output and the answer logits change. The same files record a bitwise-identical repeated clean pass (check A) and a bitwise-identical $\alpha=0$ pass (check B) for every block.

\subsection{Projection, lift, scaler, and probe}
\label{app:repro-probe}

The projection $R\in\R^{D\times512}$ is one draw of \texttt{PCG64(0).standard\_normal((D,512))} divided by $\sqrt{512}$ and stored in float32. Its entries are independent $\mathcal{N}(0,1/512)$ variates. The seed is $0$ for every checkpoint, dataset, and locus; $D$ is the family's hidden dimension. We neither normalise nor orthogonalise the columns of $R$. We normalise the lifted direction of Equation~\ref{eq:lift} once, in model space, to unit $\ell_2$ norm. The Gram matrix $R^{\top}R/D$ has full rank 512 in every one of the \cfAltdirdBlocks{} blocks that carry the displacement module, with condition number \cfAltdirdCondMin{}--\cfAltdirdCondMax{}. Table~\ref{tab:cf-altdird} reports its eigenvalue range.

The token-mean read vector $\hvec_i$ yields projected features $R^{\top}\hvec_i$. A \texttt{StandardScaler} fitted only on training rows supplies the featurewise scale $\boldsymbol s$ and mean. Every row is standardised as $(R^{\top}\hvec-\boldsymbol\mu)\oslash\max(\boldsymbol s,10^{-8})$. Each concept probe uses \texttt{LogisticRegression} with $C=1$, \texttt{solver="lbfgs"}, \texttt{max\_iter=2000}, \texttt{class\_weight=None}, and \texttt{random\_state=0}, fitted on training rows with a known label for that concept. The positive-class coefficient fixes the clinical pole of the direction. For each of the two refit seeds, we resample training units with replacement, retain every row of each sampled unit at that unit's multiplicity, and refit the scaler and six probes using the same projection. The random family, shams, and control probes remain at seed 0.

\subsection{Control constructions}
\label{app:repro-controls}

\paragraph{Random directions.} We draw one $119\times D$ standard-normal matrix with \texttt{PCG64(0)} and normalise each row to unit $\ell_2$ norm. The same 119 directions serve every question in the block.

\paragraph{Coordinate-permutation sham.} Immediately after the random draw, the same generator produces one permutation of $\{0,\dots,D{-}1\}$ per concept in protocol concept order. Concept $c$'s sham is $\hat\wvec_c$ with its coordinates rearranged by concept $c$'s own permutation. The sham preserves the direction's coordinate multiset and unit norm. Each question is written with its own sham.

\paragraph{Matched random-label probes.} Each row has a type identifier: view $\times$ sex $\times$ age decade on the chest sets, and aspect-ratio $\times$ area buckets on COCO. For control seed $k\in\{0,\dots,19\}$, we traverse the sorted type names and draw one label per type from \texttt{PCG64($k$).integers(0,2)}; if all types receive the same label, we flip the first entry. Every row of a type receives that type's label. We fit each control probe on the same training rows and known-label subset, with the same hyperparameters as the concept probe, giving it the concept probe's capacity and nuisance structure. A block fits control probes when its training rows span at least eight types. Selectivity is Equation~\ref{eq:selectivity}. Its bootstrap over calibration units keeps a draw when the real AUROC is finite and at least ten of the twenty control AUROCs are finite. We grade a cell for readability when at least 1{,}900 of the 2{,}000 draws are kept.

\subsection{Cohort roles}
\label{app:repro-cohorts}

Table~\ref{tab:cf-roles} lists each role's size and use. The roles are disjoint by patient on the chest sets and by image on COCO. We use the calibration role only to determine readability and answer capability. We fit no probe and score no write on it. The test role carries every write.

\begin{table}[h]
\centering
\scriptsize
\setlength{\tabcolsep}{5pt}
\caption{\textbf{Cohort roles.} Rows per role with independent units in parentheses, and the role's use. The units are patients on the chest sets and images on COCO.}
\label{tab:cf-roles}
\begin{tabular}{llll>{\raggedright\arraybackslash}p{0.42\textwidth}}
\toprule
Role & NIH ChestX-ray14 & CheXpert Plus & COCO & Use \\
\midrule
train & 18{,}212 (5{,}783) & 20{,}000 (20{,}000) & 20{,}000 (20{,}000) & Scaler, six concept probes, twenty control probes, the two nuisance probes of the chest sets, and the answer-direction ridge. No row of this role is scored. \\
preflight & 16 (16) & 16 (16) & 16 (16) & The seven acceptance gates, at both loci. \\
calibration & 400 (400) & 400 (400) & 400 (400) & Clean pass at $\alpha=0$: probe readability and answer capability. \\
test & 600 (600) & 600 (600) & 600 (600) & Every write matrix, dose, refit, template, locus, and addendum module. \\
valid & -- & \cfValidRows{} (\cfValidRows{}) & -- & The regrade on radiologist consensus labels (Table~\ref{tab:cf-valid}). \\
\bottomrule
\end{tabular}
\end{table}
\FloatBarrier

\subsection{Protocol record and amendments}
\label{app:repro-protocol}

Each block writes a run record containing the protocol identifier \texttt{cf-transfer-v1}, the run identifier \texttt{cf-transfer-v1:<checkpoint>:<dataset>:mayo}, the evaluation-code commit that produced the block, the pinned checkpoint and processor revisions, the library versions, the dtypes, and the resolved processing settings, including the rendered chat template. All \cfBlocks{} blocks carry the protocol identifier \texttt{cf-transfer-v1}. The protocol file registers its own code reference commit \texttt{763072d3}.

Table~\ref{tab:cf-amendments} lists additions to the protocol file after the first wave and the date of each addition. The six planned modules and the connector calibration derived from them carry no amendment date. The robustness tables in Appendix~\ref{app:cf-robustness} identify the blocks that carry each amendment module.

\begin{table}[h]
\centering
\scriptsize
\setlength{\tabcolsep}{5pt}
\caption{\textbf{Protocol amendments.} Each entry is dated by the \texttt{added} field of the registered protocol file; the first two carry the identifiers of the execution log.}
\label{tab:cf-amendments}
\begin{tabular}{ll>{\raggedright\arraybackslash}p{0.58\textwidth}}
\toprule
Amendment & Date & Change \\
\midrule
A1 & 2026-09-14 & Per-block primary template: the single-template modules score the first eligible template in the order (IY, IB), so a checkpoint whose yes/no template fails the image-free mapping gate is graded on its B-positive template. Llama~3.2 Vision 11B re-ran on all three datasets. \\
A2 & 2026-09-14 & CheXpert gains DOSE, REFIT, LOCUS, LOCUS\_CALIBRATION and PROMPT, with Effusion and Edema as its PROMPT concepts. \\
ALTDIR (A3) & 2026-09-14 & Four alternative direction constructions per concept: difference of class means, Haufe pattern, orthogonalised normal, and residualised normal. \\
EXTCOMP & 2026-09-14 & Extra clinical directions for every additional dataset label with at least 100 known positives and negatives, added to the competitor family. \\
TOKENW & 2026-09-14 & Softmax-weighted and top-quarter token weights for the six logistic directions. \\
PRECISION & 2026-09-14 & The write grid rescored with fp32 weights and forward pass, and in bf16 at batch size one. \\
ANSDIR & 2026-09-14 & The answer direction: a ridge regression of the clean answer margin on the same projected, train-scaled features. \\
ALTDIRD & 2026-09-15 & The difference-of-means and pattern vectors lifted as displacements through the minimum-norm preimage. \\
ATTR & 2026-09-15 & Three non-clinical attribute directions written with the six clinical directions as one nine-direction family. \\
ANSDIRT & 2026-09-15 & The IY-fitted answer directions written under the five held-out templates. \\
VALID & 2026-09-15 & The write grid regraded on the CheXpert validation split with radiologist consensus labels. \\
\bottomrule
\end{tabular}
\end{table}

\FloatBarrier
